\documentclass[a4paper,10pt]{article}

%%%%%%%%%%%% PREÁMBULO %%%%%%%%%%%%%%%%%%%%%

% Paquetes

\usepackage[utf8]{inputenc}
\usepackage[spanish,es-tabla]{babel}
\usepackage[T1]{fontenc}
\usepackage{graphicx}

\usepackage{listings}
\usepackage{hyperref}

\hypersetup{
    colorlinks,
    citecolor=black,
    filecolor=black,
    linkcolor=black,
    urlcolor=black
}

\usepackage{fancyhdr}
\setlength{\headheight}{15.2pt}
\pagestyle{fancy}
\fancyhead[C]{\includegraphics[width=\textwidth]{imgs/header_europa}}

% Comandos

\renewcommand{\lstlistingname}{Código}
\renewcommand{\lstlistlistingname}{Índice de fragmentos de código fuente}

% Opciones

\title{Empresas para MongoFallero}
\author{Angel Berlanas Vicente}

%%%%%%%%%%%%%%%%%%%%%%%%%%%%%%%%%%%%%%%%%%%%%%

\begin{document}
\maketitle

\begin{abstract}
Resumen de las diferentes característica y requisitos de las empresas que participan en el contrato de MongoFallero. 
\end{abstract}

\tableofcontents

\section{Requisitos Generales}

\subsection{Lenguajes}

Para la realización de las hojas de estilo que aplicarán a nuestra aplicación de MongoFallero debemos contar con el uso de \textit{SaSS} que nos permitirá realizar las hojas de estilo necesarias para adecuarnos a las empresas ganadoras y aplicar tan solo aquella que se desee en el momento del exámen (quería decir \textit{venta del producto}). De tal manera que tendremos que tener preparados los estilos haciendo uso de las posibilidades que nos da \textit{SaSS} para poder obtener esta flexibilidad.

A continuación se detallan los requisitos de diseño de cada una de lasempresas que pugnan por la aplicación, debemos de ser capaces de modificando \textit{una sola variable} en el \texttt{.scss}.

Durante la evaluación del ejercicio se pedirá que se modifique esa \texttt{variable} y que se recompilen los ficheros \texttt{css} con ayuda de \textit{SaSS}, o cualquier variación de esto:

\begin{itemize}
\item Ejecución de \texttt{ruby-sass} manualmente.
\item El uso sass con la opción de \textit{vigilancia} para estar pendiente de los cambios.
\item Configuración del módulo de \textit{nodejs} de \texttt{node-sass}.
\item Otras propuestas por el alumnado.
\end{itemize}

Una vez compilado de forma ágil y \textit{sin errores} el profesor pasará a comprobar que todos los requisitos definidos en este documento han sido cubierto para todas las empresas. Esto no significa que revise \textbf{todas}, pero si que \textit{puede hacerlo}.

Recordad que el objetivo de esta actividad es demostrar que \textit{habéis aprendido} y que tenéis las destrezas necesarias.

\includegraphics[width=\textwidth]{imgs/practice}

\subsection{Herramientas}

Debemos utilizar un editor de código y las librerias adecuadas para poner en marcha todo lo necesario.

Además deberemos tenerlo todo el GitHub, de tal manera que podamos hacer más adelante un \textit{despliegue} en \textbf{Heroku}.

El uso y dominio de la línea de comandos es un \textit{requisito}.

\section{Estructura de un proyecto SaSS}

Cuando trabajamos con herramientas más complejas, como puede ser SaSS, y teniendo en cuenta que lo más probable es que trabajemos como \textit{equipo}, es conveniente seguir una estructura concreta para organizar los diferentes ficheros.

Un posible ejemplo es el descrito en la página de SaSS, tenéis más ejemplos en los recursos que os adjunto:

\begin{itemize}
  \item \href{http://thesassway.com/beginner/how-to-structure-a-sass-project}{How-To Structure}.
  \item \href{https://scotch.io/tutorials/aesthetic-sass-1-architecture-and-style-organization}{SaSS Estético y Arquitectura}
\end{itemize}

Los nombres de los ficheros son \textit{auto-explicativos}, así que vamos a utilizarlos para nuestro proyecto de las fallas, de tal manera que quede \textit{compartimentado} todo de manera correcta y sencilla de leer \textit{y evaluar}.

\section{Empresas (themes)}

Debemos tener preparados los siguientes estilos (que deben ser generados mediante \texttt{sass}:

\subsection{ Evil Corp }

Una de las empresas más importantes del sector, pero que tiene algún problema con sus métodos de selección de personal.

\includegraphics[width=\textwidth]{imgs/evilcorp}

La \textit{Guia de Estilo} de la empresa requiere de lo siguiente.

\begin{itemize}

\item Los colores de la página han de ser complementarios del logo de la empresa.
\item Logo de EvilCorp siempre en la parte superior izquierda.
\item Logo de EvilCorp como fondo del botón de menú cuando se muestre como dispositivo móvil.
\item Los filtros de las fallas han de marcarse como seleccionados a medida que se vayan completando, de tal manera que se sepa \textit{exactamente} cúales han sido seleccionados como filtros y cuales están con su valor por defecto.
\item Las fallas con más de 4 en las puntuaciones han de aparecer con un borde diferente.
\item Los directivos de EvilCorp \textit{solo} utilizan el móvil para comprobar que las páginas están bien diseñadas, pero se comprobará también que se muestra correctamente en la versión escritorio.
  
\end{itemize}


\subsection{ Starbred Coffe }

Sin duda, la empresa más famosa de todas. 

\includegraphics[width=\textwidth]{imgs/cthulhu}

La \textit{Guia de Estilo} de la empresa requiere de lo siguiente.

\begin{itemize}

\item Los colores de la página han de ser los mismos que la página de su filial (Starbucks) \href{https://www.starbucks.es/coffee}{Ejemplo}.
\item Logo de Starbred siempre en la parte superior izquierda.
\item Menú de hamburgesa normal cuando se muestre la versión movil.
\item En la versión de escritorio los filtros aplicados han de aparecer en un cuadro superpuestos a la página aunque estemos haciendo scroll, de tal manera que sepamos en cada momento \textit{qué} estamos mostrando.
\item Las fallas con más de 4 en las puntuaciones han de aparecer con un borde y un fondo diferente.
\item Los directivos de Starbred \textit{solo} utilizan el escritorio para comprobar que las páginas están bien diseñadas, pero se comprobará también que se muestra correctamente en la versión móvil.
  
\end{itemize}



Los logos e imágenes son los que se encuentran en la carpeta adjunta.


\end{document}
